\documentclass[11pt,a4paper]{article}
\usepackage[T1]{fontenc}
\usepackage{lmodern,amsmath,amssymb,booktabs,geometry,enumitem,fancyhdr,microtype,hyperref}
\geometry{margin=22mm}
\setlength{\parindent}{0pt}
\setlength{\parskip}{6pt}
\pagestyle{fancy}
\fancyhf{}
\fancyhead[L]{Machine Learning Club}
\fancyhead[R]{Advanced 03}
\fancyfoot[C]{\thepage}
\newcommand{\problem}[1]{\subsection*{Problem #1}}
\begin{document}
{\small SEVEN HILLS INTERNATIONAL SCHOOL \hfill 2026--2027}

{\Large\bfseries Constructing and questioning an ID3 tree}\par
Worked information gain, minimality, irrelevant attributes and incomplete training.

\section*{Preliminaries}
A categorical decision tree asks a question at each internal node and returns a class at a leaf. Its \emph{size} means the number of internal decision nodes. We use the Week 3 synthetic outdoor-activity example.

\begin{center}
\begin{tabular}{cllll}
\toprule
\# & Sky & Wind & Temperature & Go?\\\midrule
1 & sunny & calm & warm & yes\\
2 & sunny & calm & cold & yes\\
3 & sunny & windy & warm & yes\\
4 & sunny & windy & cold & no\\
5 & rainy & calm & warm & no\\
6 & rainy & calm & cold & no\\
7 & rainy & windy & warm & no\\
8 & rainy & windy & cold & no\\\bottomrule
\end{tabular}
\end{center}

\section*{Worked Example A --- Choose the root}
For class proportions $p_i$, use Shannon entropy
\[
H(Y)=-\sum_i p_i\log_2 p_i,\qquad 0\log_2 0:=0.
\]
The dataset has three yes and five no labels:
\[
H(Y)=-\frac38\log_2\frac38-\frac58\log_2\frac58\approx0.9544.
\]
Among the four sunny examples, three are yes and one is no: $H(Y\mid\text{sunny})\approx0.8113$. The rainy subset is pure: $H(Y\mid\text{rainy})=0$. Thus
\[
H(Y\mid\text{sky})=\frac48(0.8113)+\frac48(0)=0.4056,
\]
\[
\operatorname{Gain}(Y,\text{sky})\approx0.9544-0.4056=\boxed{0.5488}.
\]

\problem{1 --- Compare all first questions}
\begin{enumerate}[label=(\alph*)]
\item Count yes/no labels in the calm and windy groups.
\item Calculate the weighted entropy and information gain of wind.
\item Repeat for temperature. Which root is best?
\end{enumerate}
\textbf{Selected result.} Calm has (2 yes, 2 no), windy has (1 yes, 3 no). Its weighted entropy is $(1+0.8113)/2=0.9056$ and its gain is $0.0488$. Temperature gives the same gain. Sky has the greatest gain.

\problem{2 --- Construct the complete tree}
\begin{enumerate}[label=(\alph*)]
\item Explain why rainy is immediately a no leaf.
\item Calculate gains within the sunny subset.
\item Break the resulting tie in favour of wind and draw the tree.
\item Trace all eight records; count its internal decision nodes.
\end{enumerate}
\textbf{Hint.} The sunny subset has entropy $0.8113$; its conditional entropy after splitting on wind is $0.5$. Gain is $0.3113$.
\begin{verbatim}
if sky == "rainy": return "no"
if wind == "calm": return "yes"
if temperature == "warm": return "yes"
return "no"
\end{verbatim}
This tree has three internal tests.

\problem{3 --- Minimum number of questions}
Assume every decision node tests one of sky, wind, temperature.
\begin{enumerate}[label=(\alph*)]
\item Find, for each feature, two records differing only in that feature whose labels are opposite.
\item Prove the target depends on all three features.
\item Prove that no tree with at most two internal nodes represents all eight examples correctly.
\end{enumerate}
\textbf{Selected proof.} Sky is relevant for fixed calm/cold; wind is relevant for sunny/cold; temperature is relevant for sunny/windy. Any tree with fewer than three decision nodes can test at most two distinct features; therefore it cannot represent a function that depends on all three. Problem 2 supplies a three-node construction. Three is necessary and sufficient.

\problem{4 --- An irrelevant additional attribute}
Add a fourth feature $z$, indoor-heating, which is genuinely irrelevant to the true target on the entire input domain.
\begin{enumerate}[label=(\alph*)]
\item Describe how to remove tests of $z$ from a tree without changing the function it computes.
\item Explain why this requires more than observing that $z$ seems irrelevant on a limited sample.
\end{enumerate}
\textbf{Hint.} Fix the value of $z$ consistently through the tree; the target function must remain unchanged.

\problem{5 --- One missing example}
Remove record 4 (sunny, windy, cold $\mapsto$ no) before training.
\begin{enumerate}[label=(\alph*)]
\item Construct a tree consistent with the remaining seven rows that predicts no on record 4.
\item Construct another consistent tree that predicts yes on record 4.
\item Explain why 100\% training accuracy on seven records is insufficient.
\end{enumerate}
\textbf{Hint.} The second tree may predict yes whenever the sky is sunny, because all remaining sunny examples have label yes.

\section*{Final observation}
Selecting the largest local information gain is a greedy strategy and does not by definition prove a globally minimum-size decision tree. This particular example has a separate minimality argument. A Python truth-table check is optional, but the argument must be explained in words. No special software is needed.

\section*{References}
Experience AI, \emph{Decision trees}; IOAI classical ML syllabus; the eight-row dataset and worked calculations are original to this sheet.
\end{document}
